% ৪.৪ HTML ট্যাগ, এলিমেন্ট, অ্যাট্রিবিউট ও সিনট্যাক্স
\section{HTML ট্যাগ, এলিমেন্ট, অ্যাট্রিবিউট ও সিনট্যাক্স}

\begin{learningbox}
এই টপিক শেষে শিক্ষার্থী পারবে:
\begin{itemize}
\item HTML tag, element, attribute, content ও attribute value আলাদা করতে।
\item start tag, end tag, paired tag ও empty/void tag শনাক্ত করতে।
\item valid HTML syntax ও common syntax mistake বোঝাতে।
\item ছোট HTML code বিশ্লেষণ করে প্রতিটি অংশের কাজ বলতে।
\end{itemize}
\end{learningbox}


\subsection{HTML Tag কী}
\begin{definitionbox}{Tag}
HTML document-এ content-এর ধরন বা কাজ বোঝাতে angle bracket \texttt{< >}-এর মধ্যে যে keyword লেখা হয় তাকে tag বলে। যেমন \texttt{<p>} paragraph শুরু করে এবং \texttt{</p>} paragraph শেষ করে।
\end{definitionbox}

\begin{definitionbox}{Content}
Start tag ও end tag-এর মাঝখানে যে text, image, link বা অন্য element থাকে তাকে content বলা হয়। যেমন \texttt{<h1>Welcome to ICT</h1>}-এ \texttt{Welcome to ICT} content।
\end{definitionbox}

\begin{definitionbox}{Element}
Start tag, content এবং end tag মিলিয়ে একটি HTML element তৈরি হয়। যেমন \texttt{<h1>Title</h1>} একটি heading element।
\end{definitionbox}

\begin{definitionbox}{Attribute}
Tag-এর অতিরিক্ত তথ্য বা property নির্দেশ করতে attribute ব্যবহৃত হয়। Attribute সাধারণত opening tag-এর মধ্যে লেখা হয়। যেমন \texttt{<a href="page.html">Link</a>}-এ \texttt{href} attribute।
\end{definitionbox}

\begin{definitionbox}{Attribute value}
Attribute-এর নির্দিষ্ট মানকে attribute value বলে। যেমন \texttt{<p align="center">ICT</p>}-এ \texttt{center} হলো \texttt{align} attribute-এর value।
\end{definitionbox}

\begin{keypointbox}{এক নজরে}
\begin{itemize}
\item Tag = নির্দেশনা
\item Element = start tag + content + end tag
\item Attribute = অতিরিক্ত তথ্য
\item Value = attribute-এর মান
\end{itemize}
\end{keypointbox}

\begin{center}
\begin{tabular}{|p{0.16\textwidth}|p{0.25\textwidth}|p{0.43\textwidth}|}
\hline
\textbf{Term} & \textbf{সংক্ষিপ্ত অর্থ} & \textbf{উদাহরণ} \\
\hline
Tag & angle bracket-এর মধ্যে লেখা নির্দেশনা & \texttt{<p>}, \texttt{</p>}, \texttt{<a>} \\
\hline
Element & শুরু, content ও শেষ অংশের সমষ্টি & \texttt{<p>Text</p>} \\
\hline
Attribute & tag-এর অতিরিক্ত তথ্য & \texttt{href}, \texttt{src}, \texttt{alt}, \texttt{align} \\
\hline
Attribute value & attribute-এর নির্দিষ্ট মান & \texttt{"page.html"}, \texttt{"center"} \\
\hline
Content & tag-এর মাঝের লেখা বা nested অংশ & text, image, link text, list item \\
\hline
\end{tabular}
\end{center}

\begin{figure}[H]
\centering
\resizebox{0.96\linewidth}{!}{%
\begin{tikzpicture}[
  node distance=14mm and 18mm,
  box/.style={draw, rounded corners, align=center, inner sep=4pt, text width=31mm, fill=gray!5, font=\small},
  elementbox/.style={draw, rounded corners, align=center, inner sep=4pt, text width=44mm, fill=gray!5, font=\small}
]
\node[elementbox] (element) {\textbf{HTML element}\\{\scriptsize\ttfamily <p align="center">}\\{\scriptsize\ttfamily Bangladesh}\\{\scriptsize\ttfamily </p>}};
\node[box, above left=of element] (starttag) {\textbf{Start tag}\\{\scriptsize\ttfamily <p}\\{\scriptsize\ttfamily align="center">}};
\node[box, above right=of element] (endtag) {\textbf{End tag}\\\texttt{</p>}};
\node[box, below=of element] (attr) {\textbf{Attribute}\\{\scriptsize\ttfamily align="center"}};
\draw[->, thick] (starttag) -- (element);
\draw[->, thick] (endtag) -- (element);
\draw[->, thick] (attr) -- (element);
\end{tikzpicture}
}
\caption{HTML tag, element ও attribute-এর সম্পর্ক}
\label{fig:html-tag-element-attribute}
\end{figure}

\subsection{ট্যাগের ধরন}
\begin{center}
\begin{tabular}{|p{0.22\textwidth}|p{0.30\textwidth}|p{0.34\textwidth}|}
\hline
\textbf{ধরন} & \textbf{উদাহরণ} & \textbf{কাজ} \\
\hline
Paired tag / container tag & \texttt{<p>...</p>} & content ঘিরে রাখে \\
\hline
Empty tag / void tag & \texttt{<br>}, \texttt{<img>} & end tag থাকে না; নিজেই কাজ সম্পন্ন করে \\
\hline
Attributeসহ tag & \texttt{<img src="a.jpg" alt="A">} & অতিরিক্ত তথ্য দেয় \\
\hline
\end{tabular}
\end{center}

\begin{definitionbox}{Paired tag}
যে HTML tag-এর opening tag ও closing tag দু'টি অংশ থাকে এবং মাঝখানে content থাকে, তাকে paired tag বলে। যেমন \texttt{<p>Bangladesh</p>}।
\end{definitionbox}

\begin{definitionbox}{Empty tag}
যে HTML tag-এর সাধারণত কোনো closing tag থাকে না এবং যা content ঘিরে রাখে না, তাকে empty tag বলে। আধুনিক HTML-এ এগুলোকেই প্রায়শ void element বলা হয়। যেমন \texttt{<br>}, \texttt{<hr>}, \texttt{<img>}।
\end{definitionbox}

\begin{keypointbox}{Important empty tags}
HSC MCQ-তে empty tag শনাক্ত করতে বলা হয়। মনে রাখবে:
\begin{itemize}
\item \texttt{<br>} - line break
\item \texttt{<hr>} - horizontal line
\item \texttt{<img>} - image
\item \texttt{<meta>} - metadata/charset
\item \texttt{<link>} - external resource link, যেমন CSS file
\item \texttt{<base>} - relative URL-এর base address
\end{itemize}
\end{keypointbox}

\subsection{HTML Syntax}
HTML code লেখার নিয়মকেই HTML syntax বলে। Syntax ঠিক না হলে browser অনেক সময় code-কে সঠিকভাবে বুঝতে পারে না।

\begin{definitionbox}{Syntax}
HTML document-এ tag, element, attribute ও content সঠিকভাবে লেখার নিয়মকে syntax বলে।
\end{definitionbox}

\begin{center}
\fbox{\parbox{0.92\textwidth}{
\textbf{Basic syntax pattern}\\[4pt]
\texttt{<tag>content</tag>}\\
\texttt{<tag attribute="value">content</tag>}\\
\texttt{<empty-tag>}\\[2pt]
Tag সাধারণত angle bracket-এর মধ্যে থাকে, attribute opening tag-এর মধ্যে থাকে, আর closing tag-এ slash \texttt{/} ব্যবহার করা হয়।
}}
\end{center}

\begin{keypointbox}{HTML syntax-এর প্রধান নিয়ম}
\begin{itemize}
\item Tag অবশ্যই \texttt{< >} angle bracket-এর মধ্যে লিখতে হবে।
\item Closing tag-এ keyword-এর আগে \texttt{/} slash থাকবে।
\item Attribute সাধারণত opening tag-এর মধ্যে লিখতে হবে।
\item Attribute name ও value-এর মাঝে \texttt{=} চিহ্ন থাকবে।
\item Attribute value quote-এর মধ্যে লেখা উত্তম।
\item Tag nesting সঠিক order-এ থাকতে হবে।
\item HTML tag case-insensitive, তবে lowercase tag ব্যবহার করা ভালো।
\item Empty/void tag-এর সাধারণত closing tag থাকে না।
\end{itemize}
\end{keypointbox}

\subsection{Nested HTML Element}

\begin{definitionbox}{Nested element}
একটি HTML element-এর content হিসেবে আরেকটি HTML element ব্যবহার করা হলে তাকে nested HTML element বলা হয়।
\end{definitionbox}

Nested element লিখতে হলে tag close করার order ঠিক রাখতে হয়। যে tag আগে শুরু হয়েছে, সাধারণত তার ভিতরের tag আগে close হবে, তারপর বাইরের tag close হবে।

\begin{examplebox}{Code}
\begin{verbatim}
<p>
  This is <b><i>important</i></b> text.
</p>
\end{verbatim}
\end{examplebox}

\begin{examplebox}{ব্যাখ্যা}
এখানে \texttt{<i>} tag \texttt{<b>} tag-এর ভিতরে আছে। তাই আগে \texttt{</i>} এবং পরে \texttt{</b>} লেখা হয়েছে। এটি correct nesting।
\end{examplebox}

\begin{mistakebox}{ভুল nesting}
\texttt{<b><i>ICT</b></i>} ভুল, কারণ এখানে \texttt{<i>} tag আগে close হওয়ার কথা ছিল, কিন্তু \texttt{</b>} আগে লেখা হয়েছে। তাই nesting order ভুল।
\end{mistakebox}

\subsection{Correct ও Incorrect syntax}
\begin{examplebox}{Syntax example}
\textbf{১. Paragraph tag}
\begin{verbatim}
<p>ICT is an important subject.</p>
\end{verbatim}
সঠিক, কারণ opening tag, content এবং closing tag তিনটিই আছে।

\vspace{6pt}
\textbf{২. Attributeসহ tag}
\begin{verbatim}
<p align="center">ICT</p>
\end{verbatim}
এখানে \texttt{align} হলো attribute এবং \texttt{center} হলো attribute value। বইয়ের উদাহরণে এই ধরনের লেখা board question-এ খুবই দেখা যায়।

\vspace{6pt}
\textbf{৩. Hyperlink tag}
\begin{verbatim}
<a href="next.html" target="_blank">Next Page</a>
\end{verbatim}
এখানে \texttt{href} destination page নির্দেশ করে, আর \texttt{target} link কোথায় খুলবে তা বোঝায়।

\vspace{6pt}
\textbf{৪. Image tag}
\begin{verbatim}
<img src="images/photo.jpg" alt="Photo description" width="300">
\end{verbatim}
\texttt{img} একটি empty tag; তাই এর closing tag নেই। \texttt{src} image-এর location দেখায় এবং \texttt{alt} বিকল্প text দেয়।

\vspace{6pt}
\textbf{৫. ভুল nesting}
\begin{verbatim}
<b><i>ICT</b></i>
\end{verbatim}
এটি ভুল, কারণ tag বন্ধ করার order ঠিক নেই।
\end{examplebox}

\begin{mistakebox}{সাধারণ ভুল}
\begin{itemize}
\item Tag ও element এক মনে করা।
\item Attribute-কে content মনে করা।
\item Attribute value quote ছাড়া লেখা।
\item Closing tag-এ \texttt{/} না দেওয়া।
\item Empty tag-এর closing tag লিখতে যাওয়া।
\item Tag nesting উল্টে দেওয়া।
\item \texttt{center} বা \texttt{photo.jpg}-এর মতো value-কে attribute বলে ভুল করা।
\end{itemize}
\end{mistakebox}

\subsection{Common tag examples}
\begin{examplebox}{Tag উদাহরণসমূহ}
\textbf{Document skeleton}
\begin{verbatim}
<!DOCTYPE html>
<html lang="en">
\t<head>
\t\t<meta charset="utf-8">
\t\t<title>Page Title</title>
\t</head>
\t<body>
\t\t<h1>Welcome to HTML</h1>
\t\t<p>This is a simple page.</p>
\t</body>
</html>
\end{verbatim}
\texttt{<!DOCTYPE html>} browser-কে document type জানায়; \texttt{<html>} পুরো document ধরে; \texttt{<head>} metadata/title রাখে; \texttt{<body>} দৃশ্যমান content রাখে।

\vspace{6pt}
\textbf{Headings and paragraph}
\begin{verbatim}
<h1>Main heading</h1>
<h2>Sub heading</h2>
<p>This is a paragraph.</p>
\end{verbatim}
\texttt{<h1>} সবচেয়ে বড় heading, আর \texttt{<p>} paragraph তৈরি করে।

\vspace{6pt}
\textbf{Line break and horizontal rule}
\begin{verbatim}
Line 1<br>Line 2
<hr>
\end{verbatim}
\texttt{<br>} নতুন line তৈরি করে; \texttt{<hr>} বিষয়কে আলাদা করতে একটি horizontal line দেয়।

\vspace{6pt}
\textbf{Comment}
\begin{verbatim}
<!-- This is a comment -->
\end{verbatim}
Comment ব্রাউজারে দেখা যায় না; এটি code বুঝতে সাহায্য করে।
\end{examplebox}

\begin{boardbox}{বোর্ড ফোকাস}
Tag, element, attribute ও syntax নিয়ে প্রশ্নে প্রথমে সংজ্ঞা, পরে একটি ছোট code example এবং শেষে প্রতিটি অংশের কাজ লিখলে উত্তর পরিষ্কার হয়। বিশেষ করে \texttt{<p align="center">ICT</p>} এবং \texttt{<img src="..." alt="...">} type example ভালোভাবে অনুশীলন করো।
\end{boardbox}

\begin{practicebox}
\begin{enumerate}
\item HTML tag, element ও attribute-এর পার্থক্য একটি উদাহরণসহ লেখো।
\item \texttt{<p align="center">ICT</p>} code-এ tag, attribute, value ও content শনাক্ত করো।
\item Empty tag কাকে বলে? দুটি উদাহরণ লেখো।
\item নিচের code-এ কোন syntax ভুল আছে তা শনাক্ত করো: \texttt{<h1>Welcome<h1>}
\end{enumerate}
\end{practicebox}

\begin{recapbox}
\begin{itemize}
\item Tag = angle bracket-এর মধ্যে লেখা নির্দেশনা।
\item Element = opening tag + content + closing tag।
\item Attribute = অতিরিক্ত তথ্য; value = attribute-এর মান।
\item Syntax = code লেখার নিয়ম।
\item Paired tag content ঘিরে রাখে, empty tag সাধারণত content ঘিরে রাখে না।
\end{itemize}
\end{recapbox}
